\documentclass[a4paper, 11pt]{article}
\usepackage[utf8]{inputenc}
\usepackage[T1]{fontenc}
\usepackage[italian]{babel}
\usepackage{geometry}
\usepackage{graphicx}
\usepackage{tabularx}
\usepackage{fancyhdr}
\usepackage{hyperref}

% Margini
\geometry{top=3cm, bottom=3cm, left=2.5cm, right=2.5cm}

% Intestazione e piè di pagina
\pagestyle{fancy}
\fancyhf{}
\fancyhead[L]{\small \textit{0x06tudio} -- Verbale di Riunione}
\fancyhead[R]{\small Anno Accademico 2026/2027}
\fancyfoot[C]{\thepage}
\renewcommand{\headrulewidth}{0.4pt}

\begin{document}
\thispagestyle{plain}

\begin{minipage}{0.6\textwidth}
    \vspace{-2cm}
    {\huge \textbf{0x06tudio}} \\[0.5em]
    {\large Progetto didattico di Ingegneria del Software} \\[0.3em]
    {\normalsize Università degli Studi di Padova}
\end{minipage}
\begin{minipage}{0.4\textwidth}
    \raggedleft
    \vspace{0pt}
    \includegraphics[width=6cm]{../../media/logo.png}
    \vspace{1cm}
    \small
    \href{mailto:0x06tudio@gmail.com}{0x06tudio@gmail.com}
\end{minipage}

\vspace{0.5cm}

\begin{center}
    {\LARGE \textbf{Verbale Riunione}}
\end{center}

\vspace{0.5cm}


\begin{tabularx}{\textwidth}{l X}
    \textbf{Versione:} & 0.1.0 \\
    \textbf{Data:} & 2026/10/20 \\
    \textbf{Uso:} & Esterno \\
    \textbf{Partecipanti (Presenti):} &  Nicolò Beltrame, Pietro Brocadello, Alberto Cavalletto, Catalin Chirilov, Davide Colabove, Andrea Maggio\\
    \textbf{Destinatari:} & Zucchetti Spa; Proff. Vardanega, Cardin \\
\end{tabularx}

\vspace{1cm}

\section*{Registro delle modifiche}
\begin{tabularx}{\textwidth}{|c|c|X|c|}
    \hline
    \textbf{Versione} & \textbf{Data} & \textbf{Descrizione} & \textbf{Autore} \\ \hline
    0.1.0 & 2026/10/20 & Creazione e stesura bozza iniziale & Andrea Maggio \\ \hline
\end{tabularx}

\newpage
\tableofcontents
\newpage


\section{Informazioni generali}
\begin{tabularx}{\textwidth}{l X}
    \textbf{Modalità / Luogo:} & Online / Google Meet \\
    \textbf{Data svolgimento:} & 2026/10/20 \\
    \textbf{Ora inizio:} & 14:30 \\
    \textbf{Ora fine:} & 15:30 \\
    \textbf{Durata totale:} & 1 ora \\
    \textbf{Scopo della riunione:} & Richiesta chiarimenti riguardo il capitolato C9 \\
    \textbf{Partecipanti esterni:} & Gregorio Piccoli (Zucchetti Spa)
\end{tabularx}

\newpage


\section{Ordine del giorno}
Durante la riunione sono stati affrontati i seguenti punti:
\begin{enumerate}
    \item Come gestire la context window ridotta degli LLM quando va sottoposto ad analisi un progetto di grandi dimensioni
    \item Quale harness utilizzare
    \item Come garantire un'adeguata copertura dei test automatici
    \item Copertura minima da garantire nella verifica dei casi d'uso
    \item Formati di input da preferire
    \item Proposta di un plug-in deterministico per la gestione dello stato e della dashboard
    \item Modelli locali consigliati
\end{enumerate}

\newpage


\section{Resoconto della discussione}

\subsection{Punto 1:}
Ai fini del progetto, è sufficiente anche effettuare il testing su un progetto di dimensioni ridotte, tuttavia progetti più ampi non dovrebbero creare problemi dato che questi mettono a disposizione numerosi file Markdown "di appoggio" per le IA.
Un'alternativa è quella di creare noi stessi un "progetto giocattolo" ai fini del testing.
\subsection{Punto 2:}
Claude Code è adeguato; consigliato OpenCode.
\subsection{Punto 3:}
Questione rimandata a una successiva fase del progetto; è stata proposta l'idea di costruire una skill ad-hoc.
\subsection{Punto 4:}
Considerata la natura informale delle descrizioni testuali dei casi d'uso, risulta difficile costruire dei test formali deterministici. Pertanto non è richiesta una copertura minima, e la verifica deve essere necessariamente informale.
\subsection{Punto 5:}
Preferibili i formati testuali (Markdown), lasciata al gruppo la libertà di sperimentazione con la lettura di diagrammi direttamente in formato PNG.
\subsection{Punto 6:}
Considerata appropriata la proposta di mappare ID casi d'uso con requisiti, potenzialmente anche con un database. Il plug-in potrà orchestrare poi delle chiamate all'LLM; è stato mostrato OpenSpec come esempio.
\subsection{Punto 7:}
Forniti riferimenti autorevoli per verificare l'usabilità di modelli locali sulle proprie macchine.
\newpage


\section{Approvazione Esterna}

Il presente verbale è stato condiviso con il referente aziendale in data YYYY/MM/DD ed è da ritenersi approvato e conforme a quanto discusso durante il colloquio. 

\vspace{3cm}
Firma del rappresentante:


\end{document}